%============================================================ BLOQUE 6 · Casos de uso (20 min)
\ugrsection{06}{Casos de uso}{Inicio: 1:30 · 20 min\\Cinco servicios sobre un mismo endpoint}

\ugrimageframe{imagenes/s06_diplomas.jpeg}{11.500 diplomas, un único formato}{De documentos heterogéneos a datos validados · Ilustración generada con IA}

\begin{frame}{Cinco servicios, un solo endpoint}{Todos los clientes hablan la misma API compatible con OpenAI}
	\centering\footnotesize
	\setlength{\tabcolsep}{3pt}
	\begin{tabular}{@{}R{1.9cm}R{3.2cm}R{4.8cm}R{2.3cm}@{}}
		\toprule
		\textbf{Servicio}                       & \textbf{Dónde}           & \textbf{Qué hace}                                              & \textbf{Acceso}   \\
		\midrule
		\textbf{\textcolor{rojoUGR}{MiIA Chat}} & Navegador (Open WebUI)   & Conversación, documentos, sandbox de Python, búsqueda web, voz & Keycloak (OAuth2) \\[3pt]
		\textbf{ThunderAI}                      & Extensión de Thunderbird & Reescribir, resumir, traducir y clasificar correos             & Clave virtual     \\[3pt]
		\textbf{Continue}                       & Extensión de VS Code     & Explicar, generar y refactorizar código                        & Clave virtual     \\[3pt]
		\textbf{OpenCode}                       & Terminal                 & Agente de código (bloque 5)                                    & Clave virtual     \\[3pt]
		\textbf{EPFReader}                      & Aplicación en Python     & Extracción masiva de datos                                     & Clave virtual     \\
		\bottomrule
	\end{tabular}
	\keymsg{La gente conserva sus herramientas y gana un botón más.}
\end{frame}

\begin{frame}{Caso práctico: EPFReader}{Diplomas de formación para la Carrera Horizontal del PTGAS}
	\begin{columns}[c]
		\begin{column}{0.3\textwidth}
			\bigstat{11.500+}{diplomas procesados.}
		\end{column}
		\begin{column}{0.66\textwidth}
			\footnotesize
			Campos: curso, entidad, fechas, horas, calificación.\\[4pt]
			Un proceso manual aportó la referencia de contraste.
		\end{column}
	\end{columns}
	\ugrsteps{PDF/{Diplomas de partida},
	Imágenes/{pdf2image y OpenCV},
	OCR/{Tesseract},
	LLM/{Prompt + esquema Pydantic},
	JSON/{Validado, uno por diploma}}
	\keymsg{Con un modelo de visión el OCR sobra: el modelo lee la imagen directamente.}
\end{frame}

\begin{frame}{Qué lo hizo funcionar}{Dos técnicas del bloque 3, aplicadas a gran escala}
	\begin{columns}[T]
		\begin{column}{0.5\textwidth}
			\themeblock{rojoUGR}{Salida estructurada}{%
				Un modelo Pydantic estricto define cada campo, su tipo y sus restricciones.\\
				La decodificación restringida fuerza un JSON conforme al esquema.\\
				Sin texto libre ni campos inventados.}
		\end{column}
		\begin{column}{0.5\textwidth}
			\themeblock{nA}{Visión en lugar de OCR}{%
				Llama 3.2 Vision y Qwen2.5-VL leen la imagen directamente.\\
				Mejoran con diseños heterogéneos: logotipos, sellos, firmas manuscritas, tipografías poco habituales.\\
				Justo donde el OCR clásico falla más.}
		\end{column}
	\end{columns}
	\keymsg{Documento no estructurado de entrada, registro validado de salida.}
\end{frame}

\begin{frame}{El patrón se generaliza}{Cualquier proceso que parta de documentos heterogéneos}
	{\footnotesize\color{darkgray}Candidatos:}\hspace{6pt}
	\appchip{Matrículas}\hspace{4pt}\appchip{Becas}\hspace{4pt}\appchip{Certificados}\hspace{4pt}\appchip{Facturas}\hspace{4pt}\appchip{Cuestionarios}
	\vspace{0.3cm}
	\ugrsteps{Esquema/{Define los campos},
	Extraer/{OCR o visión},
	Validar/{Reglas automáticas},
	Excepciones/{Las revisa una persona},
	Medir/{Con muestra etiquetada}}
	\vspace{8pt}\demoline{03\_salida\_estructurada.py}
	\keymsg{La receta no cambia: lo que cambia es el esquema.}
\end{frame}

\begin{frame}{Demo 03: ¿qué vamos a hacer?}{Una versión reducida del patrón de extracción de EPFReader}
	\begin{columns}[T]
		\begin{column}{0.31\textwidth}
			\ugrcard[nA]{1}{Entrada}{Un diploma inventado en texto, con errores que simulan una salida OCR.}
		\end{column}
		\begin{column}{0.31\textwidth}
			\ugrcard[nB]{2}{Proceso}{El modelo extrae los datos dos veces: sin esquema y con esquema.}
		\end{column}
		\begin{column}{0.31\textwidth}
			\ugrcard[nE]{3}{Resultado}{Pydantic comprueba los datos recibidos y podemos comparar ambas respuestas.}
		\end{column}
	\end{columns}
	{\footnotesize La demo no lee PDF ni ejecuta OCR: parte de un texto preparado.}
\end{frame}

\begin{frame}{Demo 03: comparar las dos peticiones}{El cliente comprueba ambas respuestas con el mismo modelo Pydantic}
	\centering
	\adjustbox{max width=\linewidth,max height=3.4cm}{%
		\begin{tikzpicture}[every node/.append style={font=\footnotesize}]
			\node[bbx, minimum width=3cm, minimum height=0.65cm] at (0,0) {\textbf{Programa cliente}};
			\node[rbx, minimum width=3cm, minimum height=0.65cm] at (10.4,0) {\textbf{Servidor LLM}};
			\draw[ar] (0,-0.9) -- (10.4,-0.9);
			\node[anchor=south, inner sep=1pt] at (5.2,-0.86) {1. Texto del diploma + «responde solo con JSON»};
			\draw[ar] (10.4,-1.65) -- (0,-1.65);
			\node[anchor=south, inner sep=1pt] at (5.2,-1.61) {2. Primera respuesta: el cliente intenta validarla};
			\draw[ar] (0,-2.4) -- (10.4,-2.4);
			\node[anchor=south, inner sep=1pt] at (5.2,-2.36) {3. Mismo texto + esquema JSON en \texttt{response\_format}};
			\draw[ar] (10.4,-3.15) -- (0,-3.15);
			\node[anchor=south, inner sep=1pt] at (5.2,-3.11) {4. Segunda respuesta: el cliente la valida y muestra};
		\end{tikzpicture}}
	\par\vspace{2pt}
	{\footnotesize Son dos peticiones independientes; la segunda no usa la respuesta de la primera.}
	\keymsg{La segunda petición incluye el esquema; la primera no.}
\end{frame}

\begin{frame}[fragile]{Demo 03: preparar el texto y los campos}{El ejemplo recibe texto; el esquema define qué datos esperamos}
	\begin{columns}[T]
		\begin{column}{0.5\textwidth}
			{\footnotesize\textbf{Fragmento del diploma inventado}:}
			\begin{lstlisting}[language={}]
CERT1FICA que D.ª María Ejemplo Prueba
ha superado con aprovecha-
miento el curso "Introducción a la
Administración de Sistemas Linux"
con una duración de 40 horas.
\end{lstlisting}
			{\footnotesize La entrada contiene errores y cortes de palabras.}
		\end{column}
		\begin{column}{0.48\textwidth}
			{\footnotesize\textbf{Campos definidos en Pydantic}:}
			\begin{lstlisting}
class Diploma(BaseModel):
    curso: str
    entidad_emisora: str
    alumno: str
    fecha_inicio: Optional[str]
    fecha_fin: Optional[str]
    horas: Optional[int]
    aprovechamiento: bool
\end{lstlisting}
			{\scriptsize Extracto sin las descripciones de los campos.}
		\end{column}
	\end{columns}
	\keymsg{Pydantic define el formato esperado; el texto del diploma sigue siendo la fuente de los datos.}
\end{frame}

\begin{frame}[fragile]{Demo 03: primera petición, JSON sin esquema}{Pedimos un formato en el prompt y comprobamos la respuesta}
	\begin{columns}[T]
		\begin{column}{0.58\textwidth}
			\begin{lstlisting}
r = client.chat.completions.create(
    model=MODEL, temperature=config.temperature,
    messages=[{"role": "user", "content":
        PROMPT + "\nResponde solo con JSON." + NT}],
)
bruto = r.choices[0].message.content
try:
    Diploma.model_validate_json(bruto)
    print("Válido")
except (ValidationError, ValueError):
    print("NO válido")
\end{lstlisting}
		\end{column}
		\begin{column}{0.4\textwidth}
			\footnotesize
			\begin{itemize}
				\item El servidor recibe el texto y la instrucción, pero no el esquema de \texttt{Diploma}.
				\item El modelo puede usar nombres distintos o devolver un formato incompatible.
				\item El cliente imprime el resultado, intenta validarlo y continúa aunque falle la validación.
			\end{itemize}
		\end{column}
	\end{columns}
	{\scriptsize Extracto abreviado de \texttt{03\_salida\_estructurada.py}.}
\end{frame}

\begin{frame}[fragile]{Demo 03: segunda petición, con esquema}{El cliente envía al servidor el contrato generado por Pydantic}
	\begin{columns}[T]
		\begin{column}{0.59\textwidth}
			\begin{lstlisting}
r = client.chat.completions.create(
    model=MODEL, temperature=config.temperature,
    messages=[{"role": "user",
               "content": PROMPT + NT}],
    response_format={
        "type": "json_schema",
        "json_schema": {
            "name": "Diploma",
            "schema": Diploma.model_json_schema(),
            "strict": True,
        },
    },
)
diploma = Diploma.model_validate_json(
    r.choices[0].message.content)
\end{lstlisting}
		\end{column}
		\begin{column}{0.39\textwidth}
			\footnotesize
			\begin{itemize}
				\item \texttt{model\_json\_schema()} describe los campos y tipos esperados.
				\item Si el motor lo admite, usa el esquema para restringir la generación.
				\item El cliente valida otra vez. Si esta petición o su validación falla, el ejemplo termina con error.
			\end{itemize}
		\end{column}
	\end{columns}
\end{frame}

\begin{frame}[fragile]{Demo 03: ejecutar e interpretar el resultado}{Desde \texttt{miia-demo}, con los servicios y el modelo preparados}
	\begin{lstlisting}[language={}]
docker compose run --rm lab python ejemplos/03_salida_estructurada.py
\end{lstlisting}
	\begin{columns}[T]
		\begin{column}{0.48\textwidth}
			\themeblock{nA}{Primera respuesta}{%
				Puede ser JSON válido y aun así no cumplir \texttt{Diploma}.\\
				Por ejemplo, puede usar \texttt{duracion\_horas} en vez de \texttt{horas}.}
		\end{column}
		\begin{column}{0.48\textwidth}
			\themeblock{nB}{Segunda respuesta}{%
				Si se admite el esquema y pasa la validación, se muestra el objeto como JSON.\\
				Los valores extraídos aún pueden ser incorrectos.}
		\end{column}
	\end{columns}
	{\footnotesize Las fechas son cadenas: el ejemplo no comprueba que sean fechas reales ni que sigan \texttt{AAAA-MM-DD}.}
\end{frame}

\begin{frame}{Verifica antes de fiarte}{La evaluación forma parte del sistema, no es un añadido}
	\begin{enumerate}
		\item Construye primero una \textbf{muestra etiquetada}: es tu referencia.
		\item Informa de la \textbf{precisión por campo}, no de un «parece correcto».
		\item Deriva a personas las validaciones fallidas y los casos de baja confianza.
		\item Repite la muestra al cambiar de modelo, de prompt o de motor.
	\end{enumerate}
	\keymsg{Un proceso rápido que falla en silencio es peor que uno lento que acierta.}
\end{frame}
